\documentclass[10pt]{article}

\usepackage[utf8]{inputenc}

\usepackage[T1]{fontenc}

\usepackage{amsmath}

\usepackage{amsfonts}

\usepackage{amssymb}

\usepackage[version=4]{mhchem}

\usepackage{stmaryrd}

\usepackage{mathrsfs}

\usepackage{bbold}

\usepackage{graphicx}

\usepackage[export]{adjustbox}

\graphicspath{ {./images/} }



\begin{document}

(пробегания) предметных переменных данного формального языка. В качестве формальных языков берутся языки узкого исчисления предикатов. Каждый такой язык полностью описывается множеством



\[

L=\left\{P_{0}, \ldots, P_{n}, \ldots, F_{0}, \ldots, F_{m}, \ldots\right\} \text {, }

\]



где $P_{0}, \ldots, P_{n}, \ldots$ - предикатные символы, а $F_{0}, \ldots$, $F_{m}, \ldots$ - функциональные символы, для каждого из к-рых указано число его аргументных мест. Модель $\mathfrak{M}$ (или алгебраич. система) для $L$ задаөтся нөпустым множөством $M$ и интерпретирующей функцией $I$, определенной на $L$ и сопоставляющцей $n$-местному предикатному символу $n$-местный предикат, т. е. подмножество дөкартовой степени $M^{n}$ множества $M$, a $n$-местному функциональному символу - $n$-местную функцию $M^{n} \rightarrow M$. Множество $M$ наз. П. о. (или увиверсумом) модели $\mathfrak{D}$.



Лит.: [1] К л ин и С. К., Математическая логика, пер. с англ., М., 1973; [2] К е й с ле е моделей, пер. с англ., М., 1977; [3] E р ша о Ю.. Л., П ал ю т и н Е. А., Математическая логика, М., 1979.



ПРЕДМЕТНАЯ ПЕРЕМЕННАЯ В. Н. Гришия. видиал перемениал. См. Также IIре ПРЕДМЕТНЫЙ ЯЗЫК - язЫк, являющийся предметом изучения. При формализации какой-либо содержательной теории различают два языка. Один - әто язык формализуемой теории, или П. я., задаваемый правилами построения выражений П. я. и семантическими правилами, определяющими, что обозначают его выражения или какиө они выражают суждения. Другой - это язык, на к-ром формулируются упомянутые выше синтаксические и семантич. правила. Этот язык наз. метаязыком. Обычно метаязык не формализуется. Однако и его можно формализовать, и тогда он станет П. я., для изучения к-рого требуется новый метаязык. Лит.: [1] К л и н и С. К., Введение в метаматематику,

пер. с англ., М., 1957.



ПРЕДНОРМА - то же, что полунорма.



ПРЕДПОРЯДОК, к в а 3 и п о р я д о к, п р е ду пор я доченност ь, ква ви у по р я д оч е н н о с т ь, - рефлексивное и транзитивное бинарное отношение на множестве. Если $\leq$ есть П. на множестве $M$, то отношение $a \sim b \leftrightarrow a \leq b$ и $b \leq a, a, b \in M$, является отношением эквивалентности на $M$. При этом П. $\leq$ индуцирует порядок на фактормножестве $M / \sim$. T. С. Фобанова.



ПРЕДПУЧОК на топологическом пространстве $X$ со вначениями в нек-рой категории $\mathscr{K}$ (напр., категории мвожеств, групп, модулей, колец, ...) - контравариантный бунктор $F$ из категории открытых множеств пространства $X$ и их естественных отображений включения в категорию $\mathscr{K}$. В зависимости от $\mathscr{\mathscr { K }}$ предпучок $F$ наз. П ре д пуч ком множе с тв, г р у пі, мод у л ө й, к о л е ц,... . Отвечающие вклюqениям $V \subset U$ морфизмы $F(U) \rightarrow F(V)$ наз. гомоморфизмами ограничения.



Всякий П. определяет на $X$ пучок (см. Пучков теория).



ПРЕДСКАЗУЕМАЯ $\sigma$-АЛГЕБРА - наименышая $\sigma$-алгебра $\mathscr{P}=\mathscr{P}(\mathbb{F})$ множеств из



\[

\Omega \times R_{+}=\{(\omega, t): \omega \in \Omega, t \geqslant 0\},

\]



порожденная всеми отображениями $(\omega, t) \rightarrow f(\omega, t)$ множества $\Omega \times R+$ в $R$, к-рые (для каждого фиксированного $\omega \in \Omega$ ) являются (по $t$ ) непрерывными слева и $\mathbb{F}$-согласованными с неубывающим семейством $\mathbb{F}=\left(\mathscr{F}_{t}\right)_{t \geqslant 0}$



\includegraphics[max width=\textwidth, center]{2023_07_08_5f921365a33675828d32g-1}

пространство. П. $\sigma-a$. порождается также любой из совокупностей множеств:



\begin{enumerate}

  \item $A \times\{0\}$, где $A \in \mathscr{F}_{0}$, и $\llbracket 0, \tau \rrbracket$, где $\tau$ - марковские моменты, а $\llbracket 0, \tau \rrbracket-$ стохастич. интервалы;



  \item $A \times\{0\}$, где $A \in \mathscr{F}_{0}$, и $A \times(s, t]$, где $s<t$ и $A \in \mathscr{F}_{s}$. Между опциональными алгебрами и П. б-а. имеет место соотношение $\mathscr{P}(\mathbb{F}) \equiv \mathbb{O}(\mathbb{F})$.



\end{enumerate}



Лum.: [1] Д е л л а ш е р п К., Емкости и случайные про-



ПРЕДСКАЗУЕМЫЙ СЛУЧАИНЫЙ ПРОЦЕСС стохастическшй процесс $X=\left(X_{t}(\omega), \mathcal{F}_{t}\right)$, являющийся измеримым (как отображение $(\omega, t) \rightarrow X(\omega, t)=X_{t}(\omega)$ )



\begin{center}

\includegraphics[max width=\textwidth]{2023_07_08_5f921365a33675828d32g-1(1)}

\end{center}



ПРЁСТА виМОСТИ МАТРИЦ проблема, заключающаяся в том, чтобы выяснить, можно ли указать такой единый общий метод (алгоритм) к-рый по провзвольной системе $U, U_{1}, \ldots, U_{q}$ целочисленных матриц позволял бы за конечное число шагов ответить на вопрос, представима ли матрица $U$ через остальные матрицы $U_{1}, \ldots, U_{q}$ с помощью операции умножения. Наибольший интерес представляөт случай, когда матрицы $U, U_{1}, \ldots, U_{q}$ являются квадратными и имеют один и тот же порядок. Так сформулированная П. м. П. наз. о б щ е й. Фиксируя матрицы $U_{1}, \ldots, U_{q}$ и оставляя матрицу $U$ переменной, получают т. н. ч а с т н ы е П. м. П. Алгоритм, решающий общую П. м. п., решал бы и все частные проблемы, так что для установления неразрешимости общей П. М. І. достаточно указать хотя бы одну неразрешимую частную П. м. П.



П. М. П.- одна из первых алгоритмических проблем алгебраич. характера, неразрешимость к-рых была установлена. Первоначально А. А. Марковым было показано (см. [1], [2]), что для любого $n \geqslant 6$ может быть построена система, состоящая из 91 матрицы порядка $n$, такая, что соответствующая ей частная П. м. П. будет неразрешимой, т. е. будет невозможен алгоритм (понимаемый в точном смысле этого слова), распознающий по произвольной матрице порядка $n$, представима ли она через матрицы данной системы. В дальнейшем (см. [3]) число матриц в системе было уменьшено до 23 и было показано, что за счет надлежащего усложнения конструкции системы (включая увеличевие числа членов) условие $n \geqslant 6$ может быть ослаблено до $n \geqslant 4$. Для любого $n \geqslant 6$ строится конкретная система, состоящая из 12 матриц порядка $n$, с неразрешимой частной П.м. П. (см. [4]). Путем подходящей фиксации $U$ и варьирования $U_{1}, \ldots, U_{q}$ доказана неразрешимость общей П. М. п. для $n=3$ (см. [5]).



Лит.: [1] M a p K о в А. А., "Докл. АH CCCP", 1951, т. 78, № 6, с. 1089-92; [2] е го ж е, Теория алгорифмов, М.- Ј., 1954 (Tp. MaTeM. HH-ra AH CCCP, T. 42); (3] e ro He, "Z. Math.

Logik und Grundl. Math.", 1958, Bd 4, S. 157-68; [4] H aг о р в ы й H. M., "VI Bcecoюзн. конференция по матем. логй кe», Тб., 1982, c. 124; [5] P a t e rs o n M. S., "Studies in

appl. math.", 1970, v. $49, \mathcal{N}_{1,}$ p. 105-07. H. M. Hazopныü. appl. math.", 1970, v. 49, № ФУ. (или контравариантный) функтор $F$ из нек-рой категории $\mathfrak{R}$ в категорию множеств $\Im$, изоморфный одному из основных теоретико-множественных функторов:



\[

H_{A}(X)=H(A, X): \Re \rightarrow \Im .

\]



Функтор $F: \Re \rightarrow S$ представим тогда и только тогда, когда найдутся такие объект $A \in \mathrm{Ob} \mathfrak{R}$ и элемент $a \in F(A)$, что для каждого элемента $x \in F(X), X \in \mathrm{Ob} \mathfrak{\Re}$, существует единственный морфизм $\alpha: A \rightarrow X$, для к-рого $x=a F(\alpha)$. Объект $A$ называется п р е д с т а в л я ющ и м ф у к т о р $F$; он определен однозначно с точностью до изоморфизма.



В категории множеств тождественный функтор представим: представляющим объектом служит одноточечное множество. Функтор взятия нек-рой декартовой степени также представим: представляющим объектом служит множество, мощность к-рого равна этой степени. В произвольной категории произведение П. ф. $F_{i}$ с представляющими объектами $A_{i}, \quad i \in I$, представимо тогда и только тогда, когда в этой категории существует копроизведение объектов $A_{i}$ Всякий ковариантный П. Ф. перестановочен с пределами, т. е. непрерывен.





\end{document}